\begin{lemma}
Let $1<p<\infty$, let $n\ge0$ be an integer, let $\Omega\subseteq\mathbb{C}$ and
let $H\colon\mathbb{C}\to\mathbb{C}$ belong to $W^{1,2}_{\mathrm{loc}}(\Omega)$
in the sense of \noderef{MemWLocOn} and solve equation (18) on $\Omega$ in the
sense of \noderef{SolvesEq18}.  Put
\[
  K=\max\Bigl\{p-1,\ \frac1{p-1}\Bigr\} .
\]
Then $H$ is $K$-quasiregular on $\Omega$ in the sense of
\noderef{QuasiregularOn}.

This is the paper's assertion at line 239, ``Since $H$ is $K$-quasiregular,
$K=\max\{p-1,1/(p-1)\}$'', which the paper states without proof and which the
estimates (26) and (27) then use.  Because $K$ here is determined by $p$ alone,
a constant of the form $C(n,p,K)$ obtained from it is a constant $C(n,p)$;
this is what allows Theorem 2 to carry no quasiregularity hypothesis.
\end{lemma}
\begin{proof}
By \noderef{QuasiregularOn} we must verify three things: $1\le K$, that
$H\in W^{1,2}_{\mathrm{loc}}(\Omega)$, and that there are weak Wirtinger
derivatives $a=\partial H/\partial\xi$, $b=\partial H/\partial\bar\xi$ of $H$
on $\Omega$, in the sense of \noderef{IsWeakWirtingerOn}, with
\begin{equation}\tag{11}
  |b(\xi)|\le\frac{K-1}{K+1}\,|a(\xi)|
  \qquad\text{for almost every }\xi\in\Omega .
\end{equation}
The second is the hypothesis.  We prove the first and third.  Note that the
derivatives in the hypothesis and in the conclusion are the same objects: by
\noderef{SolvesEq18} the system (18) is an identity between weak Wirtinger
derivatives of $H$ on $\Omega$, and by \noderef{QuasiregularOn} so is the
distortion inequality, so no passage between two notions of derivative is
needed here.

\emph{Step 1: $p-1>0$, and $1\le K$.}  Since $p>1$ we have $p-1>0$, so
$1/(p-1)$ is defined and positive.  If $p\ge2$ then $p-1\ge1$, so
$K\ge p-1\ge1$.  If $p<2$ then $0<p-1<1$, so $1/(p-1)>1$ and $K\ge1/(p-1)>1$.
In both cases $1\le K$.

\emph{Step 2: $\dfrac{K-1}{K+1}=\dfrac{|p-2|}{p}$.}  We compute the two cases
separately.

If $p\ge2$ then $p-1\ge1$, hence $1/(p-1)\le1\le p-1$, so $K=p-1$.  Therefore
$K-1=p-2$ and $K+1=p$, and since $p-2\ge0$,
\[
  \frac{K-1}{K+1}=\frac{p-2}{p}=\frac{|p-2|}{p}.
\]

If $1<p<2$ then $0<p-1<1$, hence $p-1<1<1/(p-1)$, so $K=1/(p-1)$.  Multiplying
numerator and denominator by the positive number $p-1$,
\[
  \frac{K-1}{K+1}=\frac{\frac1{p-1}-1}{\frac1{p-1}+1}
  =\frac{1-(p-1)}{1+(p-1)}=\frac{2-p}{p}=\frac{|p-2|}{p},
\]
the last equality because $p-2<0$.  In both cases
\begin{equation}\tag{1}
  \frac{K-1}{K+1}=\frac{|p-2|}{p}.
\end{equation}

\emph{Step 3: the estimate off the origin.}  By \noderef{SolvesEq18} there are
weak Wirtinger derivatives $a,b$ of $H$ on $\Omega$, in the sense of
\noderef{IsWeakWirtingerOn}, and a Lebesgue null set $Z$ such that
\begin{equation}\tag{18}
  b(\xi)=\Bigl(\frac12-\frac1p\Bigr)
    \Bigl[\frac{\xi}{\bar\xi}a(\xi)+\frac{\bar\xi^{\,n}}{\xi^{n}}
      \overline{a(\xi)}\Bigr]
\end{equation}
for every $\xi\notin Z$ with $\xi\in\Omega$ and $\xi\neq0$.  Fix this pair
$a,b$, and fix such a $\xi$.  Since $\xi\neq0$ we have $\bar\xi\neq0$ and
$|\bar\xi|=|\xi|\neq0$, so the two bracketed factors are defined and
\[
  \Bigl|\frac{\xi}{\bar\xi}\Bigr|=\frac{|\xi|}{|\bar\xi|}=1,\qquad
  \Bigl|\frac{\bar\xi^{\,n}}{\xi^{n}}\Bigr|
    =\frac{|\bar\xi|^{n}}{|\xi|^{n}}=1 .
\]
Also $\bigl|\overline{a(\xi)}\bigr|=|a(\xi)|$.  Hence by (18), the triangle
inequality and multiplicativity of the modulus,
\[
  |b(\xi)|
  \le\Bigl|\frac12-\frac1p\Bigr|
    \Bigl(\Bigl|\frac{\xi}{\bar\xi}\Bigr||a(\xi)|
      +\Bigl|\frac{\bar\xi^{\,n}}{\xi^{n}}\Bigr|
        \bigl|\overline{a(\xi)}\bigr|\Bigr)
  =2\Bigl|\frac12-\frac1p\Bigr|\,|a(\xi)| .
\]
Finally, since $p>1>0$,
\[
  \Bigl|\frac12-\frac1p\Bigr|=\Bigl|\frac{p-2}{2p}\Bigr|=\frac{|p-2|}{2p},
\]
so $2\bigl|\tfrac12-\tfrac1p\bigr|=|p-2|/p$ and, by (1),
\begin{equation}\tag{2}
  |b(\xi)|\le\frac{|p-2|}{p}\,|a(\xi)|
  =\frac{K-1}{K+1}\,|a(\xi)| .
\end{equation}

\emph{Step 4: removing the restriction $\xi\neq0$.}  Step 3 establishes (2) at
every $\xi\in\Omega-Z$ with $\xi\neq0$.  The singleton $\{0\}$ is a Lebesgue
null set in $\mathbb{C}\cong\mathbb{R}^2$, and the union of two null sets is
null; outside $Z\cup\{0\}$ every $\xi\in\Omega$ satisfies $\xi\neq0$ and hence
(2).  So (2) holds for almost every $\xi\in\Omega$, which is (11) for the pair
$a,b$ of Step 3.

Together with Step 1 and the hypothesis $H\in W^{1,2}_{\mathrm{loc}}(\Omega)$,
the three clauses of \noderef{QuasiregularOn} hold for the constant $K$, the
pair exhibited in its clause (b) being the pair $a,b$ of Step 3; this is the
assertion.  Note that no property of $n$ was used beyond $n\ge0$: the factor
$\bar\xi^{\,n}/\xi^{n}$ has modulus $1$ for every such $n$.
\end{proof}
